% =====================================================================
%  Front matter: executive summary + keywords
%  Booklet: 250-350 words, one continuous block. Written LAST.
%  Owner: shared (front matter)
% =====================================================================
\chapter*{תקציר מנהלים}
\addcontentsline{toc}{chapter}{תקציר מנהלים}
\markboth{תקציר מנהלים}{}

דוח מכין זה מציג תוכנית השקה בת \ltr{6} חודשים בישראל עבור \en{AURA}, מותג
טינט שפתיים בהשראת מזון, לנשים בנות \ltr{18--28}. הרקע הוא שתי מגמות עולמיות:
מוצרי השפתיים הפכו למנוע צמיחה באיפור, ומוצרי יופי מתגלים יותר ויותר ברשתות
החברתיות. שוק הקוסמטיקה בישראל מוערך ב-\ltr{\$408.53} מיליון בשנת \ltr{2026},
וקהל היעד המיידי כולל כ-\ltr{211} אלף סטודנטיות. הבעיה היא שמותג חדש מתחיל
מאפס מודעות, בשוק שבו מתחרי הפרימיום עולים \ltr{99--175} ש"ח ונרכשים מרחוק,
ומוצר שנבחר לפי גוון, מרקם וניחוח קשה למכור בלי התנסות. לכן הוגדרו חמישה יעדי
\en{SMART}, ובהם מודעות של \ltr{20\%} לפחות בקמפוסי הפיילוט ומכירות של \ltr{255}
אלף ש"ח בחודשים \ltr{4--6}, ומפרט של \ltr{11} דרישות עם ערכי יעד ושיטות אימות.
סקירת \ltr{12} מאמרים שפיטים מראה שגירויים חושיים, שמות מפתיעים, משפיענים
אמינים ודוגמיות חינם אפקטיביים כל אחד בנפרד, אך שילובם בהשקה של מותג קוסמטיקה
לא נבדק; מאמר הדגל מספק מודל למדידת השפעת הדוגמיות דרך הרחבה, האצה
וקניבליזציה, וההשקה תמדוד בהתאם כמה מהקונות בדוכנים הן קונות חדשות באמת. חקר
השוק מזהה חלל \en{Masstige} סביב מחיר של \ltr{69} ש"ח, בין מוצרי הפארם הזולים
לבין מותגי הפרימיום המיובאים, ומעריך את השוק הזמין בפלח הסטודנטיות בכ-\ltr{25.55}
מיליון ש"ח בשנה; יעד ההשקה הוא נתח קטן ושמרני ממנו. מבין שלוש חלופות, נבחרה
במטריצת החלטה משוקללת החלופה ההיברידית (ציון \ltr{8.10}, לעומת \ltr{7.40}
ו-\ltr{6.45}): אתר \en{D2C}, מרקטפלייס ישראלי, \ltr{20--30} משפיעניות קטנות
ושלושה ברי טעימות בקמפוסים, שבהם הקונות טועמות את המזון ומנסות את הגוון.
תוכנית העבודה מחלקת את ההשקה לשש חבילות עבודה ושש אבני דרך בתקציב של \ltr{150}
אלף ש"ח, עם שער \en{Go/No-Go} בסוף החודש השלישי, שבו נבדקות הודעות השיווק
למשרד הבריאות. ניתוח הסיכונים מזהה שבעה סיכונים, ותוכנית ההפחתה מורידה את
החשיפה המרבית מציון \ltr{16} ל-\ltr{6}. המסקנה: התוכנית מוכנה לביצוע, והצעד
הראשון הוא הגשת ההודעות למשרד הבריאות.

\keywords{השקת מותג \sep ביוטי בהשראת מזון \sep שיווק משפיענים \sep דוגמיות \sep שיווק חושי \sep אומניצ'אנל}
